% SGA 3, Expose V, Loop 1 English reconstruction.
% Scope: printed pp. 253--255, from Lemma 1.2 through the terminology at
% the end of section 1; section 2 is excluded.
% Authority: Polo--Gille Exp5-13oct24.pdf, component-local pp. 3--5
% (combined-reader pp. 265--267). The PDF controls text, formulas,
% numbering, notes, and the temporary Loop-1 diagram images.
% Jacob Reinhold's CC BY 4.0 Markdown is comparison/drafting material only.

% Source page: Exp5 component-local p. 3; combined-reader p. 265;
% printed p. 253.
\medskip
\phantomsection\label{sga3:exp5:lemma:1.2}
\noindent\textbf{引理 1.2.}---
\textit{设 \(d_0,d_1:X_1\rightrightarrows X_0\) 是
\((\mathbf{Sch})\) 中的图, \((Y,p)\) 是它在
\((\mathbf{Esp.An})\) 中的余核.
\begin{enumerate}
\item[(i)]\phantomsection\label{lem:sga3-V-1.2-i}%
若 \(Y\) 是概形且 \(p\) 是概形态射,
则 \((Y,p)\) 是 \((\mathbf{Sch})\) 中 \((d_0,d_1)\) 的余核.
\item[(ii)]\phantomsection\label{lem:sga3-V-1.2-ii}%
假设 \(X_0\) 的每个点都有一个满足下述性质的饱和开邻域 \(U\).
令 \(\widetilde d_0\) 与 \(\widetilde d_1\) 为
\(d_0\) 与 \(d_1\) 在 \(d_0^{-1}(U)=d_1^{-1}(U)\) 上的限制,
令 \((Q,q)\) 为 \((\widetilde d_0,\widetilde d_1)\)
在 \((\mathbf{Esp.An})\) 中的余核.
若 \(Q\) 是概形且 \(q\) 是概形态射, 则 \((Y,p)\)
是 \((\mathbf{Sch})\) 中 \((d_0,d_1)\) 的余核.
\end{enumerate}}

\hyperref[lem:sga3-V-1.2-i]{(i)} 在
\hyperref[proof:sga3-V-4.1-c]{\S4(c)} 中证明;
由于证明很短, 这里再写一遍.
设 \(f:X_0\to Z\) 是满足 \(fd_0=fd_1\) 的概形态射.
由假设, 存在唯一的赋环空间态射 \(r:Y\to Z\) 使 \(f=rp\).
只须证明, 对每个 \(y\in Y\), \(r\) 诱导的同态
\(\mathscr O_{r(y)}\to\mathscr O_y\) 都是局部同态.
这是因为 \(p\) 满射, 因而存在某个 \(x\) 使 \(y=p(x)\),
并且 \(f\) 诱导的同态 \(\mathscr O_{f(x)}\to\mathscr O_x\)
是局部同态.

\hyperref[lem:sga3-V-1.2-ii]{(ii)} 由
\hyperref[lem:sga3-V-1.2-i]{(i)} 和
\hyperref[lem:sga3-V-1.1]{前一引理} 得到.

\medskip
\phantomsection\label{par:sga3-V-1-c}%
\noindent\textbf{c)} 本篇研究当双箭头 \((d_0,d_1)\) 处于更丰富的结构中时,
\(\operatorname{Coker}(d_0,d_1)\) 的存在性.
更确切地说, 设
\[
  X_2=X_1\times_{d_1,d_0}X_1
\]
为下图的纤维积
% Native Loop-2 reconstruction; printed p. 253.
\begin{equation}\SGAInlineEquationTarget{sga3:v:lemma:1-2:equation:star}
  \resizebox{0.20\textwidth}{!}{%
    \SGAThreeNativeInput{figures/exp5/p253_star_fiber_product.png}}
  \tag{*}\label{fig:sga3-exp5-star-fiber-product}
\end{equation}
设 \(d'_0,d'_2:X_2\to X_1\) 为两个典范投影.
于是, 按照定义, 方形
% Native Loop-2 reconstruction; printed p. 253.
\begin{equation}\SGAInlineEquationTarget{sga3:v:lemma:1-2:equation:0}
  \resizebox{0.21\textwidth}{!}{%
    \SGAThreeNativeInput{figures/exp5/p253_square_0.png}}
  \tag{0}\label{fig:sga3-exp5-square-0}
\end{equation}
是 Cartesian 方形.

此外, 假设给定第三个箭头 \(d'_1:X_2\to X_1\).
若对 \(\mathscr C\) 的每个对象 \(T\), \(X_1(T)\)
都是某个群胚 \(X_\ast(T)\) 的箭头集, 其对象集为 \(X_0(T)\),
源映射为 \(d_1(T)\), 目标映射为 \(d_0(T)\), 复合映射为
\(d'_1(T)\), 则称
\[
  (d_0,d_1:X_1\rightrightarrows X_0,d'_1)
\]
为 \(\mathscr C\)-\emph{群胚}. 如常, 作等同
\[
 (X_1\times_{d_1,d_0}X_1)(T)
 \simeq X_1(T)\times_{d_1(T),d_0(T)}X_1(T),
\]
并回顾一下, 群胚就是每个箭头均可逆的范畴.
\footnote{\label{note:sga3-V-5}编者注: 此时 \(X_2(T)\)
是可复合箭头对 \((f_2,f_1)\) 的集合, 即满足
\(d_0(f_1)=d_1(f_2)\) 的箭头对;
映射 \(d'_0,d'_1,d'_2\) 分别将 \((f_2,f_1)\) 映到
\(f_2,f_2\circ f_1,f_1\).}

% Source page: Exp5 component-local p. 4; combined-reader p. 266;
% printed p. 254.
若 \(\varphi\) 是群胚 \(X_\ast(T)\) 的箭头,
则映射 \(f\mapsto\varphi\circ f\) 是一个双射:
从目标等于 \(\varphi\) 的源的箭头 \(f\) 的集合,
映到与 \(\varphi\) 有相同目标的箭头的集合.
等价地, 这表示方形
% Native Loop-2 reconstruction; printed p. 254.
\begin{equation}\SGAInlineEquationTarget{sga3:v:lemma:1-2:equation:1}
  \resizebox{0.19\textwidth}{!}{%
    \SGAThreeNativeInput{figures/exp5/p254_square_1.png}}
  \tag{1}\label{fig:sga3-exp5-square-1}
\end{equation}
是 Cartesian 方形.

类似地, \(g\mapsto g\circ\varphi\) 是一个双射:
从群胚 \(X_\ast(T)\) 中源等于 \(\varphi\) 的目标的箭头 \(g\) 的集合,
映到与 \(\varphi\) 有相同源的箭头的集合.
等价地, 方形
% Native Loop-2 reconstruction; printed p. 254.
\begin{equation}\SGAInlineEquationTarget{sga3:v:lemma:1-2:equation:2}
  \resizebox{0.19\textwidth}{!}{%
    \SGAThreeNativeInput{figures/exp5/p254_square_2.png}}
  \tag{2}\label{fig:sga3-exp5-square-2}
\end{equation}
是 Cartesian 方形.

另一方面, 令 \(s:X_0\to X_1\) 为 \(\mathscr C\) 中的唯一箭头,
使得对每个 \(T\), 映射 \(s(T):X_0(T)\to X_1(T)\)
将 \(X_\ast(T)\) 的每个对象映到其恒等箭头.
\footnote{\label{note:sga3-V-6}编者注: \(T\mapsto s(T)\)
定义了 \(\operatorname{Hom}(h_{X_0},h_{X_1})\) 的一个元素,
而由 Yoneda 引理, 后者等于 \(\operatorname{Hom}(X_0,X_1)\).}
箭头 \(s\) 满足
\begin{align}
  d_1s&=\operatorname{id}_{X_0},
    \tag{3}\label{eq:sga3-exp5-3}\\
  d_0s&=\operatorname{id}_{X_0}.
    \tag{3 bis}\label{eq:sga3-exp5-3bis}
\end{align}
最后, 复合映射 \(d'_1(T)\) 的结合性等价于下图交换
% Native Loop-2 reconstruction; printed p. 254.
\begin{equation}\SGAInlineEquationTarget{sga3:v:lemma:1-2:equation:4}
  \resizebox{0.72\textwidth}{!}{%
    \SGAThreeNativeInput{figures/exp5/p254_associativity_4.png}}
  \tag{4}\label{fig:sga3-exp5-associativity-4}
\end{equation}

% Source page: Exp5 component-local p. 5; combined-reader p. 267;
% printed p. 255.
反之, 条件 \hyperref[sga3:v:lemma:1-2:equation:1]{(1)},
\hyperref[sga3:v:lemma:1-2:equation:2]{(2)} 和
\hyperref[sga3:v:lemma:1-2:equation:4]{(4)}, 再加上存在满足
\hyperref[eq:sga3-exp5-3]{(3)} 的箭头 \(s\), 蕴含
\((X_1\rightrightarrows X_0,d'_1)\) 是 \(\mathscr C\)-群胚.
条件 \hyperref[eq:sga3-exp5-3]{(3)} 无关紧要:
它只是保证对每个 \(T\in\mathscr C\),
\(d_1(T):X_1(T)\to X_0(T)\) 为满射.
在下文中, 主要使用 Cartesian 方形
\hyperref[sga3:v:lemma:1-2:equation:0]{(0)},
\hyperref[sga3:v:lemma:1-2:equation:1]{(1)} 和
\hyperref[sga3:v:lemma:1-2:equation:2]{(2)}, 可归纳于下图
% Native Loop-2 reconstruction; printed p. 255.
\begin{equation}\SGAInlineEquationTarget{sga3:v:lemma:1-2:equation:0-1-2}
  \resizebox{0.34\textwidth}{!}{%
    \SGAThreeNativeInput{figures/exp5/p255_diagram_012.png}}
  \tag{0,1,2}\label{fig:sga3-exp5-diagram-012}
\end{equation}
此图左侧的两个方形 (即方形
\hyperref[sga3:v:lemma:1-2:equation:0]{(0)} 和
\hyperref[sga3:v:lemma:1-2:equation:2]{(2)}) 是 Cartesian 方形;
第一行正合; 并将 \(X_2\) 等同于纤维积 \(X_1\times_{d_0,d_0}X_1\).

我们只间接地使用结合性, 例如用来保证存在满足
\hyperref[eq:sga3-exp5-3]{(3)} 和
\hyperref[eq:sga3-exp5-3bis]{(3 bis)} 的箭头 \(s\),
或保证存在箭头
\begin{equation}
  \sigma:X_1\longrightarrow X_1,
  \qquad d_0\sigma=d_1,\quad d_1\sigma=d_0.
  \tag{\ensuremath{\dagger}}\label{eq:sga3-exp5-dagger}
\end{equation}
这里选择 \(\sigma\), 使 \(\sigma(T):X_1(T)\to X_1(T)\)
将 \(X_\ast(T)\) 的每个箭头映到其逆箭头.
\footnote{\label{note:sga3-V-7}编者注: Yoneda 引理表明,
\(\sigma\) 是 \(X_1\) 的一个对合自同构; 例如,
在 \hyperref[par:sga3-V-3-e]{3(e)} 和
\hyperref[thm:sga3-V-4.1]{定理~4.1} 中将用到这一点.}

作为术语滥用, 当
\hyperref[sga3:v:lemma:1-2:equation:0]{(0)},
\hyperref[sga3:v:lemma:1-2:equation:1]{(1)} 和
\hyperref[sga3:v:lemma:1-2:equation:2]{(2)} 都是 Cartesian 方形,
\hyperref[sga3:v:lemma:1-2:equation:4]{(4)} 交换,
并且存在满足 \hyperref[eq:sga3-exp5-3]{(3)} 的 \(s\) 时,
有时将图
% Native Loop-2 reconstruction; printed p. 255.
\begin{equation*}
  \resizebox{0.40\textwidth}{!}{%
    \SGAThreeNativeInput{figures/exp5/p255_groupoid_triple_arrow.png}}
  \phantomsection\label{fig:sga3-exp5-groupoid-triple-arrow}
\end{equation*}
称为 \(\mathscr C\)-群胚.
因此, \(X_2\) 可以是 \hyperref[sga3:v:lemma:1-2:equation:star]{(*)}
的``一个''纤维积, 而非 \hyperref[sga3:v:lemma:1-2:equation:star]{(*)}
的``那个''纤维积.
\footnote{\label{note:sga3-V-8}编者注: 参见下面的
\hyperref[par:sga3-V-2-a]{例~2(a)}.}

\medskip
\phantomsection\label{sga3:exp5:terminology:c-groupoid}
\noindent\textit{术语.} 对于 \(\mathscr C\)-群胚 \(X_\ast\),
我们也称之为\emph{以 \(X_0\) 为底的群胚 \(X_\ast\)},
或 \emph{\(X_0\) 中的预等价关系 \(X_\ast\)}.

\section*{2. \(\mathscr C\)-群胚的例子}
\phantomsection\label{sec:sga3-V-2}
\addcontentsline{toc}{section}{2. \texorpdfstring{\(\mathscr C\)}{C}-群胚的例子}

\phantomsection\label{par:sga3-V-2-a}%
\noindent\textbf{a)} 设 \(X\) 是 \(\mathscr C\) 的对象,
\(G\) 是在 \(X\) 上左作用的 \(\mathscr C\)-群.
用 \(d_0:G\times X\to X\) 表示定义 \(G\) 在 \(X\) 上的作用的箭头,
用 \(d_1:G\times X\to X\) 表示到第二个因子的投影,
并用 \(\mu:G\times G\to G\) 表示定义
% Translation statistics: source_lines=221; output_lines=217; inline_math=110; display_math=3; equation=7; equation_star=1; align=1; footnotes=4.
% Review: complete assigned component; pending_TODO=0; all math symbols, labels, references, diagram paths and environments checked; formula prose translated.
